% !TEX root = main.tex

\subsection{Benchmark suites and experimental setup}
We evaluate our solver on a representative selection of SMT-COMP benchmarks, including all instances from SMT-COMP 2024\footnote{\url{https://zenodo.org/communities/smt-lib/}} and the full set of problems accepted for SMT-COMP 2025\footnote{(available at \url{https://github.com/SMT-LIB/benchmark-submission})}.  From the combined pool of roughly 100\,000 \verb|QF_S| and \verb|QF_SLIA| problems, we randomly sample 2\,000 instances proportional to each track’s share of the total, so that our test set reflects the underlying distribution of problem types.
All experiments were run on a 24 $\times$ 3.1 GHz-core server with 185 GB RAM, invoking each solver in a single thread with a 60s timeout and no memory limit. 


\subsection{Performance evaluation}

The experiments are conducted on the following solvers: cvc5~\cite{cvc5}, Z3~\cite{z3}, Z3-alpha~\cite{z3alpha}, Z3-Noodler~\cite{noodler}, OSTRICH1~\cite{ostrich}, and different engines of OSTRICH2\footnote{\url{https://zenodo.org/records/15378807}} and OSTRICH2 running the engines in a time-sliced portfolio using the flag \texttt{-portfolio=strings}.
For the RCP engine we evaluate two complementary configurations: one mode that combines forward and backward propagation without Nielsen splitting (using the flags \texttt{+F+B-N}), and another mode (using the flags \texttt{-F+B+N}) that disables forward propagation in order to emphasize Nielsen's equation decomposition. 
In practice, combining forward propagation and splitting yields limited synergy: the extra word equations generated by Nielsen splitting create significant propagation overhead, so omitting forward propagation often improves overall solving performance.  

\begin{figure}[htbp]
  
\centering
\resizebox{\linewidth}{!}{%
    \begin{tabular}{lccccc}
        \toprule
                        & CVC5 & Z3   & Z3-alpha & Z3-Noodler    & OSTRICH1 \\
        \midrule
        \textsc{sat}    & 1135 & 1074 & 1077     & \textbf{1192} & 755      \\
        \textsc{unsat}  & 722  & 726  & 701      & \textbf{744}           & 645      \\
        \midrule
        unknown/timeout & 144  & 200  & 222      & \textbf{64}            & 600      \\
        \midrule
        Solved          & 1856 & 1800 & 1778     & \textbf{1936}          & 1400     \\
        \bottomrule
    \end{tabular}%
}
  \caption{The experiment result of cvc5, Z3, Z3-alpha, Z3-Noodler and OSTRICH1 on the SMT-COMP'24/25 benchmarks.}
  \label{fig:experiment_res_z3_cvc5}
\end{figure}

\begin{figure}[htbp]
  
\centering
\resizebox{\linewidth}{!}{%
  \begin{tabular}{lccccc}
    \toprule
    \textbf{}       & OSTRICH2      & ADT  & CE   & RCP-F+B+N & RCP+F+B-N \\
    \midrule
    \textsc{sat}    & \textbf{1177}          & 586  & 960  & 1015      & 1001      \\
    \textsc{unsat}  & \textbf{760}  & 560  & 678  & 744       & 748       \\
    \midrule
    unknown/timeout & \textbf{63}   & 854  & 362  & 241       & 231       \\
    \midrule
    Solved          & \textbf{1937} & 1156 & 1638 & 1759      & 1769      \\
    \bottomrule
  \end{tabular}%
}

  \caption{The experiment result of different engines of OSTRICH2 on the SMT-COMP'24/25 benchmarks.The flags \texttt{+F+B-N} enable forwards/backwards propagation and disable the NielsenSplitter. The flags \texttt{-F+B+N} disable forwards propagation and enable backwards propagation and the NielsenSplitter.}
  \label{fig:experiment_res_ostrich}
\end{figure}


\noindent Figure~\ref{fig:experiment_res_z3_cvc5} shows that among the off‐the‐shelf solvers, Z3‐Noodler leads with 1\,936 solved problems, followed by cvc5 (1\,856), Z3 (1\,800), Z3-alpha (1\,778), and OSTRICH1 (1\,400).  Figure~\ref{fig:experiment_res_ostrich} shows that both RCP configurations (1\,759 and 1\,769 solves) and CE‐Str (1\,638) represent clear upgrades over OSTRICH1, whereas the ADT‐Str engine (1\,156) trails—largely due to its simple algorithm.
%
\noindent Importantly, when we combine all four engines in a time‐sliced portfolio, OSTRICH2 solves 1\,937 problems, putting it on par with Z3‐Noodler and ahead of the rest of the field.  We further observe that Z3‐Noodler excels on the \textsc{sat} instances, while OSTRICH2 performs particularly well on the \textsc{unsat} cases.

Finally, although Z3‐Noodler does not support \texttt{replace\_all}, only 2.4\% of the SMT-COMP benchmarks use this string function, so its absence has only a modest effect on the overall rankings. 
We remark that, since our evaluation uses a stratified random sample of 2\,000 instances drawn from the full pool of roughly 100\,000 \verb|QF_S| and \verb|QF_SLIA| benchmarks, there is inevitably some statistical variance in the exact per‐solver ranking; nonetheless, the fact that OSTRICH2 solves essentially as many problems as the leader demonstrates its overall competitiveness.  

